\section*{September 9, 2026: Single-family homes, 2008--2018}
\textit{Drafted by Codex.}

At Jacob's request, the quarter-mile descriptive packet and annual event studies
now have a single-family version. Corrected assessor classes 202--210, 234,
278, and 295 retain houses and townhouses; class-211 apartment buildings are
excluded. Noah's 30/49 school definitions, spatial exclusions, transaction
weighting, and 2022-dollar adjustment remain fixed. The original citywide
price-per-square-foot cutoffs are retained, so this change isolates property
selection. These sales are an exact subset of the earlier packet.

The descriptive sample contains 5,645 transactions: 1,461 treated and 4,184
control. Missing characteristics remain allowed. The treated mean falls from
\$288,603 in 2016 to \$247,263 in 2018, while the control mean rises from
\$191,686 to \$218,465. The treated decline therefore persists after excluding
apartment buildings. It does not establish falling prices for fixed homes:
which homes and neighborhoods supply sales can still change.

\includegraphics[width=0.95\textwidth]{../input/single_family_price_trends.png}

\small The ten-page packet is \texttt{data\_overview\_single\_family.pdf} in
\texttt{tasks/audits/home\_data\_overview/output}. Its composition page shows
building size and age. Its count stages begin with corrected single-family
classes, already restricted to records with one building card.
\normalsize
\clearpage

\subsection*{Annual event studies on the single-family sample}
All four models use 5,643 sales with complete characteristics at 75 sites.
They retain school-site and year fixed effects, 2012 as the reference year,
and school-site clustered 95\% intervals. Apartment count is omitted because
it is constant zero. Annual sample counts and mean prices match the packet's
complete-characteristics summaries.

The 2018 estimate with property controls is \$14,513 (95\% interval:
$-\$54,328$ to \$83,355), versus \$57,561 in the earlier sample including
apartment buildings. The single-family log estimate is $-0.161$ (interval:
$-0.431$ to $0.108$), versus $0.005$ in the earlier sample. Both single-family
intervals include zero. These changes reflect both the restriction in property
type and its consequences for neighborhood sales weights; they do not isolate
a causal role for apartment composition. Pre-trend and anticipation concerns
remain.

\includegraphics[width=0.95\textwidth]{../input/single_family_event_studies.png}

\small Reproduction: run \texttt{make ../output/data\_overview\_single\_family.pdf}
in \texttt{tasks/audits/home\_data\_overview/code}, and
\texttt{make ../output/event\_studies\_single\_family.pdf} in
\texttt{tasks/audits/home\_prices\_twfe/code}. Saved summary datasets have
standard reports. This is a September 9 working-tree analysis; the broader
property samples remain available.
\normalsize
